\section{Exercises}\label{chap:measure:ex}

\paragraph{On Section~\ref{ch:distances}.}
\begin{exercises}
\item Prove Proposition~\ref{prop:raw} for a positive measure $\nu$ on a general measurable space that does not charge the atoms of $U$.
\item Show that the total variation distance between the unit point masses at $a$ and $b$ is $2$ for $a\ne b$, and deduce that total variation cannot compare two atomic measures with slightly displaced atoms.
\item Compute the $L^1$ distance between the cumulative functions of a Cauchy kernel of width $\varepsilon$ and a unit step on $[-L,L]$, derive the $\varepsilon\log(1/\varepsilon)$ rate, and compare with the formula for a single step in Section~\ref{sec:ch9:mass}.
\item Derive the first order formula of Proposition~\ref{prop:mass} for the sinh clock of Section~\ref{ch:ggc}, where it reads $-(\varepsilon/\pi)f'/f$ with $f(\theta)=\sin(\pi\theta)/(\pi\theta)$, and compare it at $\Theta=m+\frac12$ with the exact mass $m+\pi^{-1}\arctan(\varepsilon/\Theta)$ of Section~\ref{sec:ch3:bench}; explain why the difference is of order $\varepsilon^3$.
\item Explain why the two point Richardson combination removes the linear term in $\varepsilon$ and the three point combination the quadratic term, and reproduce the heat trace errors $2.6\times10^{-7}$ and $1.7\times10^{-10}$ from the values of $W_\varepsilon(0.01)$ given in Section~\ref{sec:ch7:cumulative}.
\end{exercises}

\paragraph{On Section~\ref{ch:offline}.}
\begin{exercises}
\item Verify the expansion $\frac12(P_{\varepsilon+\delta}+P_{\varepsilon-\delta})-P_\varepsilon=\frac{\delta^2}2\partial_a^2P_a|_{a=\varepsilon}+O(\delta^4)$, show that the remainder has $L^1$ norm $O(\delta^4/\varepsilon^4)$, and compute the $L^1$ norm of $\partial_a^2P_a$ at $a=\varepsilon$.
\item Show that for $\delta>\varepsilon$ the kernel $P_{\varepsilon+\delta}-P_{\delta-\varepsilon}$ is negative exactly on $|x|<\sqrt{\delta^2-\varepsilon^2}$.
\item Compute the depth $2\varepsilon/(\pi(\delta^2-\varepsilon^2))$ of the negative dip of an off line pair at its centre, that is $\varepsilon/(\pi(\delta^2-\varepsilon^2))$ per displaced zero, compare it with the contribution $\varepsilon/(\pi(\varepsilon^2+\delta_j^2))\approx\varepsilon/(\pi\delta_j^2)$ of a line zero at distance $\delta_j$, and derive the local visibility condition of Section~\ref{ch:finite}.
\item Given the matched floor $\tilde d_{\rm match}(\varepsilon)$, derive the smallest displacement $\delta$ of a pair that keeps the count and would be detected, and evaluate it at $\varepsilon=0.05$ using Table~\ref{tab:tv}.
\end{exercises}

\paragraph{On Section~\ref{ch:turing}.}
\begin{exercises}
\item Show that $I(\varepsilon)=\int_0^T\log|\xi(\frac12+\varepsilon+it)|\,dt$ is continuous for $\varepsilon\ge0$, that its derivative jumps by $2\pi$ when the line $\re s=\frac12+\varepsilon$ crosses a zero, and deduce the second identity of Proposition~\ref{prop:littlewood}.
\item Prove Proposition~\ref{prop:zerodensity} in detail, paying attention to the orientation of each side of the rectangle.
\item\label{ex:ch9:pair} Show that the kernels $P_{\varepsilon+\delta}$ and $-P_{\delta-\varepsilon}$ of an off line pair at distance $\delta>\varepsilon$ from the line have total mass zero, and compute to leading order the mass that they and their mirror images at $-\gamma$ put on a window $(0,T]$ containing the ordinate $\gamma$. Explain how the factor two in Proposition~\ref{prop:zerodensity} accounts for the pair.
\item Using the first order formula for $\tau_\varepsilon(T)$ in Section~\ref{sec:ch9:count}, predict $\tau_{0.05}(200)$ from $\Xi'(200)/\Xi(200)=-1.2506879$ and compare with Table~\ref{tab:zerodensity}; estimate the coefficient of the $\varepsilon^3$ term from the rows $\varepsilon=0.01$ and $0.02$.
\item\label{ex:ch9:littlewood} Show that Littlewood's lemma follows from Proposition~\ref{prop:littlewood} and Proposition~\ref{prop:zerodensity} by integrating in $\varepsilon$.
\end{exercises}

\paragraph{On Section~\ref{ch:sin2more}.}
\begin{exercises}
\item Show that $\frac1\pi\int_0^\Theta(1-\cos\theta x)\,d\theta=\frac1\pi(\Theta-\sin(\Theta x)/x)$ and, for $\alpha>1$, derive the Dirichlet series for $\arg\zeta(\alpha+i\Theta)$ from the integrated sine squared lemma.
\item Implement Selberg's formula at $s=\frac34+i\theta$ with $x=100$ and check that adding the zero sum over the first hundred zeros reduces the error to the size of the omitted zeros.
\item Estimate, from the single zero law, the cutoff $x$ needed for six correct digits of $v_{0.55}(\theta)$ at an ordinate $\theta=\gamma$ from the primes, and compare with the cost of the P\'olya route at the same point.
\item Show that $g(0)-g(x)\ge0$ for every real positive definite function $g$, and that $g(0)-g(x)$ is conditionally negative definite.
\end{exercises}
